\section{Results}
\label{sec:results}

\subsection{Warm-start}

% Generated by scripts/export_results_table.py -- do not edit by hand.
\begin{table}[t]
  \centering
  \caption{Warm-start ranking quality on the held-out test partition ($K=10$).
  Hyperparameters were fixed before the test partition was read once.}
  \label{tab:results}
  \begin{tabular}{lrrrrr}
    \toprule
    Model & NDCG@10 & Recall@10 & MAP@10 & Hit rate & Coverage \\
    \midrule
    ID-only MF & \textbf{0.1120} & 0.2084 & 0.0827 & 0.2084 & 0.8950 \\
    Popularity & 0.0788 & 0.1508 & 0.0572 & 0.1508 & 0.1000 \\
    Hybrid WARP & 0.0655 & 0.1286 & 0.0466 & 0.1286 & 0.8900 \\
    Random & 0.0364 & 0.0687 & 0.0265 & 0.0687 & 1.0000 \\
    \bottomrule
  \end{tabular}
\end{table}


On the warm-start test partition, the ID-only MF baseline leads on raw
ranking quality (Table~\ref{tab:results}). This mirrors the honest negative
result reported by the sibling BPR project: a model with fewer structural
constraints has an easier time fitting a warm population it has already
seen reviews from, because every identity it is scored on had its own row to
learn. The hybrid model pays a complexity tax here -- it has to route all of
its signal through a comparatively small vocabulary of shared features
rather than a dedicated parameter per identity -- and that tax is not fully
recovered by the identity feature it is also given.

\subsection{Cold-start}

% Generated by scripts/export_results_table.py -- do not edit by hand.
\begin{table}[t]
  \centering
  \caption{Cold-start retrieval quality ($K=10$). Users/items in this
  slice were withheld from training entirely; the ID-only baseline and the
  popularity baseline can only fall back to a popularity prior for them,
  while the hybrid model scores them through shared category, brand and
  ingredient features.}
  \label{tab:cold-start}
  \begin{tabular}{llrrrr}
    \toprule
    Model & Slice & NDCG@10 & Recall@10 & MAP@10 & Hit rate \\
    \midrule
    Hybrid WARP & New user, new item & 0.0282 & 0.0758 & 0.0126 & 0.0909 \\
    Random & New user, new item & 0.0236 & 0.0389 & 0.0098 & 0.1212 \\
    Popularity & New user, new item & 0.0000 & 0.0000 & 0.0000 & 0.0000 \\
    ID-only MF & New user, new item & 0.0000 & 0.0000 & 0.0000 & 0.0000 \\
    Random & Known user, new item & 0.0308 & 0.0570 & 0.0160 & 0.1084 \\
    Hybrid WARP & Known user, new item & 0.0129 & 0.0328 & 0.0047 & 0.0623 \\
    Popularity & Known user, new item & 0.0000 & 0.0000 & 0.0000 & 0.0000 \\
    ID-only MF & Known user, new item & 0.0000 & 0.0000 & 0.0000 & 0.0000 \\
    Popularity & New user, known item & 0.2710 & 0.1378 & 0.1333 & 0.9250 \\
    ID-only MF & New user, known item & 0.2661 & 0.1268 & 0.1363 & 0.8500 \\
    Hybrid WARP & New user, known item & 0.1956 & 0.0970 & 0.0888 & 0.7750 \\
    Random & New user, known item & 0.0642 & 0.0438 & 0.0197 & 0.5250 \\
    \bottomrule
  \end{tabular}
\end{table}


Table~\ref{tab:cold-start} is the result this project exists to produce. On
the \emph{new user, known item} slice, every model (including the hybrid
one) relies on an item-level bias/popularity signal, so all three
non-random baselines perform comparably; the hybrid model's demographic
feature sum (skin type and tone) is a weaker but genuinely non-trivial
signal on its own, clearing the random floor by a wide margin.

On the \emph{known user, new item} and \emph{new user, new item} slices, the
comparison is structural rather than a matter of degree: the popularity and
ID-only baselines both report \textbf{exactly zero} NDCG, Recall and MAP.
This is not a close loss -- a withheld item has a training count of
precisely zero, so its popularity-bias term and its identity row are both
untrained, and neither baseline has any other signal to fall back on for it.
The hybrid model scores strictly above zero on both slices, because its item
representation routes through brand, category and ingredient-tag embeddings
learned from every \emph{other} item that shares those tags -- exactly the
mechanism Equation~\ref{eq:feature-sum} was built to provide.

\subsection{Honest caveats}

The hybrid model's absolute recall on the cold-item slices is modest, and
on this synthetic fixture's small catalogue ($\sim$200 products) a uniform
random baseline is occasionally competitive with it in raw hit-rate, simply
because the ratio of $K$ to catalogue size is already favourable to chance.
We report this rather than tuning the random seed until it looks better: the
qualitative claim this project supports is that the hybrid model is
\emph{capable of saying anything at all} about a withheld entity, which the
two comparison baselines are structurally incapable of, not that it is
dramatically better than chance on every metric at this dataset scale. A
larger catalogue, more training epochs, or a sharper per-feature affinity
signal would be the natural next experiments to widen the absolute margin
(Section~\ref{sec:conclusion}).
